\section{The permanent and the rook predictor}\label{app:torres}
\label{app:rook}
\paragraph{The G1 reduction [D].}
When comparing $\operatorname{per}M$ with $\operatorname{per}M_e$, both the size and the doubly
stochastic scaling change. We split the error $g_T$ into two quantities that can be computed at
the same scaling. The incidence vector $a_e$ has a one in the row and another in the column of
$e$. Let $\tau_e=p_e(1-p_e)/(1-2p_e)$ and $S_T=\sum_e\tau_e a_ea_e^\top$. The product of its
positive eigenvalues is $\operatorname{pdet}S_T$. With sums over the allowed cells, the
normalization at both endpoints of an increment is
\[
\begin{aligned}
 H_T(p)&=\sum_e[-p_e\log(p_e/w_e)+(1-p_e)\log(1-p_e)],\\
 L_T(M)&=H_T(p)-\tfrac12\sum_e\log\frac{1-2p_e}{(1-p_e)^2}
         -\tfrac12\log\operatorname{pdet} S_T+\tfrac12\log(4\pi k).
\end{aligned}
\]
Scaling the rows by $u_i$ and the columns by $v_j$ increases $H_T$ and the logarithm of the
permanent by the same sums of logarithms. The other terms depend on $p$, which is invariant under
this scaling. Thus $g_T(M)=g_T(P)$. For a doubly stochastic matrix $D$ of size $K$, define
\[
 b_D=D-\mathbb U_K,\qquad
 \mathcal E(D)=\log\frac{K^K\operatorname{per}D}{K!}
             +\tfrac12\log\det(I-b_D^\top b_D).
\]
The logarithm of the permanent cancels exactly in $\Delta=g_T(M)-\mathcal E(P)$. Two calculations
remain. For $\mathcal E(P)$ we need the first correction to the leading determinant in the
permanent approximation. For $\Delta$ it suffices to expand the explicit terms of the
normalization.

\paragraph{The static estimate $(L^*)$ [D].}
The leading determinant for the diffuse class with a gap is derived in~\cite{dw}. McCullagh's
expansion~\cite{mcc} identifies formal corrections, and Li's estimate~\cite{li} treats
perturbations with sublinear line norm. Here we need the error after extracting the coefficient of
$1/K$ to be $O(K^{-2})$ with the same constant throughout the class. To prove this, we collect the
minors of the centered kernel in the polynomial
\[
 \mathcal R_b(s)=\sum_{\ell=0}^K s^\ell
         \sum_{|I|=|V|=\ell}\operatorname{per} b[I,V],\qquad b=b_D .
\]
Expanding $\operatorname{per}(\mathbb U_K+b)$ and using $\mathbb E T^{-\ell}=(K-\ell)!/K!$ for
$T\sim\Gamma(K+1,1)$ gives $K^K\operatorname{per}D/K!=\mathbb E\,\mathcal R_b(K/T)$. For
$b=U\Sigma V^\top$ and a standard complex Gaussian $Z$, Wick's rule identifies $\mathcal R_b(s)$
with $\mathbb E\prod_a(1+w_a)$, where $w=\sqrt{s}(U\Sigma^{1/2}Z,V\Sigma^{1/2}\overline Z)$. The
linear terms sum to zero. Absorbing $\exp(-\sum_a w_a^2/2)$ into the density gives
\[
 \mathcal R_b(s)=h(s)\widetilde{\mathbb E}\prod_a (1+w_a)e^{-w_a+w_a^2/2},\qquad
 h(s)=\det(I-s^2b^\top b)^{-1/2}.
\]
For real $s$ near one, $C=\widetilde{\mathbb E}ww^\top=sB_0(I+sB_0)^{-1}$, where $B_0$ has zero
diagonal blocks and off-diagonal blocks $b,b^\top$. Each row or column contributes a factor
beginning with $1+w^3/3-w^4/4$. At first order there is one vertex of degree four or two of degree
three. There are three pairings at the quartic vertex. In the cubic case, six pairings give three
edges between the vertices, and nine give a loop at each vertex. The expression with sums over all
indices is
\[
 \frac{\mathfrak f(s)}K=-\frac34\sum_a C_{aa}^2+
 \sum_{a,a'}\left(\tfrac12 C_{aa}C_{aa'}C_{a'a'}+\tfrac13 C_{aa'}^3\right).
\]
Take $\delta=\gamma/4$. For $|s-1|\le\delta$ the gap gives propagator entries of order $O(K^{-1})$
and bounded absolute row sums. The product requires distinct indices. The correction to the sum
over all pairs is $-\frac56\sum_a C_{aa}^3=O(K^{-2})$, since $a$ runs over the $2K$ rows and
columns. This correction is retained at the next order.

Let $\mathcal W_b(s)=\mathcal R_b(s)/h(s)$ and let $f_m(s)$ be the coefficient of $t^m$ in
$\mathcal W_b(st)$. Let $f_m^{(j)}$ denote its part of excess $j$ (edges minus vertices). For
$N=\lceil\kappa\sqrt K\rceil$, subtracting the excess-one class summed over all orders leaves
\[
 \mathcal W_b(s)-1-\frac{\mathfrak f(s)}K
 =-\frac56\sum_a C_{aa}^3+\sum_{j=2}^{N-1}\sum_{m\le N}f_m^{(j)}
   +\sum_{m>N}f_m-\sum_{m>N}f_m^{(1)} .
\]
For $2\le j<N$, each truncated class is bounded by $2j(j\Lambda/K)^j$ with uniform $\Lambda$.
Their sum is $O(K^{-2})$, since the ratio of successive bounds is below $1/2$ up to this cut-off
for large $K$. The tail of the interpolant is bounded by $C e^{C_1\sqrt K}r^{-N}$ with $r>1$. For
the excess-one tail, the diagrams have two or three edges. Distributing the order $m$ among $d$
edges gives $\binom{m-1}{d-1}$ terms with decay $\theta^{m-d}$, where
$\theta=(1+\delta)\max(1-\gamma,1/2)<1$. Thus the tail is bounded by $C
K^{-1}(N+1)^2\theta^{N-2}$. With $\kappa\log r>C_1$, both tails decrease exponentially in $\sqrt
K$, uniformly for $s\in[1-\delta,1+\delta]$ and over the entire class. Outside $|K/T-1|\le\delta$
we use $|\mathcal R_b(s)|\le\mathbb E(1+s\sum_l\sigma_l|Z_l|^2/K)^K$. A shift in the gamma density
absorbs the power. Its tails and $\sum_l\sigma_l\le\sqrt{K(M_0-1)}$ give $O(\exp(-cK+C\sqrt K))$.
Since $\mathbb E(K/T)=1$ and $\mathbb E(K/T-1)^2=1/(K-1)$, the expansion with third and fourth
moments of order $O(K^{-2})$ gives
\[
 \mathcal E(D)=\mathfrak f_{\rm e}(KD)/K+O(K^{-2}),\qquad
 \mathfrak f_{\rm e}=\mathfrak f(1)+h''(1)/(2h(1)).
\]
The term involving $h''(1)$ accounts for the variation of the determinant when averaging over
$K/T$. To obtain the coefficient of the error $g_T$, we still need to add the contribution of
$\Delta$ to this expression.
\subsection{The profile coefficient and its stability}
To add the two corrections in G1, it is convenient to express them in terms of the same profile
$\nu=kP$. When composing kernels we use the uniform row and column measures, so that each sum
over a fiber is divided by the board size. Write $B=\nu-1$, $A_R=BB^\top$, $A_C=B^\top B$,
$T_X=A_X(I-A_X)^{-1}$, $U=T_RB=BT_C$, and $V=B+U$. For $X=R,C$, let
$j_X(t)=(T_X-A_X)(t,t)$, $t_X(t)=T_X(t,t)$, and $j_0=\int j_R=\int j_C$.
The cubes in double integrals are powers of the kernel values. In this notation, the functional
that enters (T) is
\begin{align*}
 \mathfrak F={}&\tfrac12j_0^2+\operatorname{tr}(T_R^2)-\tfrac32 j_0
    -\tfrac34\sum_X\int j_X^2+\langle j_R,V j_C\rangle\\
 &-\tfrac12\sum_X\langle j_X,T_Xj_X\rangle-\tfrac13\sum_X\iint T_X^3
    +2\iint BU^2+\tfrac23\iint U^3,\qquad
 \Phi=\mathfrak F-\tfrac1{12}.
\end{align*}
Let $\bar e$ be the projection of $\nu$ in $L^2(\nu)$ onto the orthogonal complement of the
additive row and column functions. Set $m_R=\int\nu^2\,dy$, $m_C=\int\nu^2\,dx$. The
deterministic part of the coefficient is
\begin{align*}
 \delta_1={}&\tfrac1{12}+\tfrac56\iint\nu^3+\tfrac12\iint\nu\bar e^2
 -\iint V\nu\bar e-\tfrac12\iint\nu^2\\
 &+\tfrac12\iint\nu^2(t_R+t_C)-\iint\nu^2V-\tfrac14\sum_X\int m_X^2 .
\end{align*}
To obtain $\delta_1$ from the normalization, we relate the stationary point $p$ to the scaling
$P$. There are potentials $\alpha,\beta$ such that
$P_{ij}=\exp(\alpha_i+\beta_j)p_{ij}(1-p_{ij})$. Multiplying the logarithm by $P_{ij}$ and
summing eliminates the gauge. Substituting in the exact expression for $\Delta$ and expanding
the entropy and the pseudodeterminant gives $\Delta=\delta_1/k+O(k^{-2})$.
The expansion uses diffusion to bound the power sums and the gap to bound the inverses.
To verify the coefficient, use $V B^\top=T_R$, $B^\top V=T_C$, $T_X A_X=T_X-A_X$, and the fact
that $B$ has zero marginals. If $g=\iint B^2$ and $\gamma_X(t)=A_X(t,t)$, the projection
above takes the form
\[
 \bar e=B+g-\rho_R-\rho_C,\quad
 \rho_R=(I+T_R)\gamma_R-V\gamma_C,\quad
 \rho_C=(I+T_C)\gamma_C-V^\top\gamma_R .
\]
Substituting this expression and $t_X=\gamma_X+j_X$ into $\delta_1$ and the static formula
gives $\mathfrak f_{\rm e}+\delta_1=\Phi$ on each grid. Hence $R=R_{\rm est}+R_{\rm det}$, with
$R_{\rm est}=\mathcal E(P)-\mathfrak f_{\rm e}/k$ and $R_{\rm det}=\Delta-\delta_1/k$.
Both are $O(k^{-2})$, proving (T). The ratio of two permanents involves the difference between
their errors. The decomposition lets us estimate this difference separately for $R_{\rm est}$ and
$R_{\rm det}$.

\paragraph{Static stability $(L^{*+})$ [D].}
For $(L^*)$ it sufficed to bound all diagrams of excess at least two. Once the board changes, we
must determine which of their contributions survive in the minor. Retain the diagrams of excess
two and three, with the correction for coincident indices $-\frac56\sum C_{aa}^3$. The same
bounds leave a tail of order $O(k^{-4})$. The gamma average and the logarithm give
$R_{\rm est}=f_2/k^2+f_3/k^3+O(k^{-4})$, with $f_2,f_3$ uniformly bounded. On passing to a
minor, the mass defect per fiber is $O(k^{-1})$. The potentials that correct it have supremum
norm $O(k^{-1})$. Removing a single square concentrates the defect on two fibers. The gauge-fixed
inverse has bounded supremum norm and off-diagonal entries of order $O(k^{-1})$. The potentials
are $O(k^{-1})$ on the two fibers and $O(k^{-2})$ elsewhere. The perturbation argument preserves
the class with bounds $2M_0,\gamma/2$. Expanding the difference of a product of propagators marks
the edge that changed. The remaining sums are bounded over a spanning tree using the absolute row
norm and the entrywise bound. Restricting the index to the affected fiber gains a factor of
$k^{-1}$ for a localized deletion. The coefficients $f_2,f_3$ change by $O(k^{-1})$ on passing
to the minor and by $O(k^{-2})$ under deletion. Accounting also for the change in the denominators
gives orders $k^{-3}$ and $k^{-4}$ for the respective increments of $R_{\rm est}$.

\subsection{A marked square and the determinant increment [D]}
It remains to estimate the increment of $R_{\rm det}$. We can measure the effect of $e=(r,c)$
on $L_T$ by varying its weight or by comparing the board with the minor left after occupying it.
For an allowed $e$, let $\pi^{\rm der}_e=\partial_{\log w_e}L_T=p_e(1+\hat c_e)$,
$\varepsilon_e=p_e/P_e-1$, and $r_e=L_T(M_e)-L_T(M)-\log(\pi^{\rm der}_e/w_e)$.
Theorem A for rooks gives $|\hat c_e|\le C/k$ and $|r_e|\le C/k^2$. Thus $r_e$ measures the
discrepancy between the two operations through second order. Using the scaling of the minor
$P'=\operatorname{diag}(x')P^{(e)}\operatorname{diag}(y')$, where $P^{(e)}$ is the submatrix
of $P$, gives the exact identity
\begin{align*}
 R_{\rm det}(M_e)-R_{\rm det}(M)&=-r_e-\mathsf E_e,\\
 \mathsf E_e={}&\log(1+\varepsilon_e)+\log(1+\hat c_e)
          +\sum_{i\ne r}\log x'_i+\sum_{l\ne c}\log y'_l\\
 &+(k-1)\log(1-1/k)
  +\tfrac12\log\frac{\det(I-b_{P'}^\top b_{P'})}{\det(I-b_P^\top b_P)}\\
 &+\frac{\delta_1((k-1)P')}{k-1}-\frac{\delta_1(kP)}k .
\end{align*}
Indeed, $\Delta=-L_T(P)+\log(k!/k^k)-\frac12\log\det(I-b_P^\top b_P)$.
The covariance of $L_T$ under scaling gives the two sums of potentials. The sign of the size term
comes from the ratio of the two factorials. To prove $(A^+)$ we still need to gain one power of
$k$ in the signed sum $r_e+\mathsf E_e$. The exact expression for $L_T$ retains the quadratic
displacement of the stationary point, usually denoted by $\zeta$. Bounding the eleven second-order
remainders in absolute value would lose their cancellation.

\paragraph{Independent parameters.}
To retain the signed sum, we place both boards in the same family. We vary separately the mass of
row $r$, the mass of column $c$, the inverse size, and the weight of square $e$. Fix $W=kP$ and
take each row and column as a type. With $u=1/k$, let $A_-,B_-$ be the uniform probability
measures that omit the distinguished row and column, respectively. Set
\[
  A_s=(1-s)A_-+s\delta_r,\qquad B_t=(1-t)B_-+t\delta_c,\qquad h=1/J .
\]
Multiply only $W_{rc}$ by $v$. With respect to these measures, the scaled and stationary kernels
satisfy $N=W^v\exp(a_i+b_l)$ and $q(1-hq)=W^v\exp(z_i+y_l)$ with unit marginals.
With the gauge fixed, the inverse of the derivative is uniformly bounded in supremum norm.
The gap controls the quadratic norm, and the integral part maps the quadratic norm to the
supremum norm because $W\le M_0$. Perturbing the measures costs $O(|s-u|+|t-u|)$, and the
marked weight costs $O(\max(|s|,|t|))$ for $|v|<2$. Iteration with the inverse at the center
converges on fixed disks for $h,s,t$ depending only on the class and gives the holomorphic branch.
The determinants of the centered kernels $V$ that occur take the form
\[
 D(V)=\sum_{\ell\ge0}\frac{(-1)^\ell}{(\ell!)^2}
  \int\det[V(i_a,l_b)]_{a,b\le\ell}^{\,2}\,
               dA_s^{\otimes\ell}dB_t^{\otimes\ell}.
\]
Cauchy--Binet and Hadamard give the bound $\ell^\ell C^\ell/(\ell!)^2$. The gap and the
Hilbert--Schmidt norm at the center give a positive lower bound for the determinant there.
The series and Cauchy's estimates let us take its logarithm on uniformly smaller disks. In the
expression for $\Delta$ we replace $\log(J!/J^J)-\frac12\log(4\pi J)$ with
$-1/h-\frac12\log2+h/12$ and keep the Stirling remainder $\eta_J=O(J^{-3})$ separate.
For the resulting expression $D_*$, the product $h^2D_*$ is bounded and holomorphic on these disks.
For positive rational masses, replicating each type $J A_s(i)$ or $J B_t(l)$ times preserves
the operator on types and adds only zero singular values. Applying the deterministic part of G1
to these replicas gives $D_*=\delta_1/J+O(J^{-2})$. Thus the first three coefficients of $h^2D_*$
vanish and the coefficient of degree three is $\delta_1$. By continuity and the holomorphic
identity theorem, Cauchy's estimates give $R_{\rm det}=\eta_J+h^2D_2(h,s,t,v)$, with $D_2$ and
its first derivatives bounded on smaller disks. On a face of zero mass, only the visible part is
replicated. Its gauge drops out of the sums of potentials, and its own $\delta_1$ is identified.

Passing to the minor replaces $(h,s,t)=(u,u,u)$ by $(u/(1-u),0,0)$. The change in $h$ is
$O(u^2)$ and the change in each mass is $O(u)$, giving $\Delta_e R_{\rm det}=O(u^3)$.
The identity above proves $(A^+)$ with the stated sign. On the face $s=0$ or $t=0$, the
integrated expressions are independent of $v$. Thus $\partial_vD_2=stH$ with $H$ bounded by
Cauchy's estimates. At $s=t=u$ the derivative of $R_{\rm det}$ with respect to $v$ is $O(u^4)$
all the way to $v=0$. Integration over $[0,1]$ gives the deletion estimate in $(D^+)$. Adding
the increments from $(L^{*+})$ gives $(T^+)$ for the total remainder $R$. Thus we can compare
the approximation for the same pair before and after passing to a further minor that retains it.

\subsection{From the two-square marginal to the square}
\label{app:pair}
To reintroduce a diagonal restriction, we need the probability that the two rooks occupy its
endpoints. We estimate it first with the ratio of predictors $\pi^T$ and then with the local
expression $\pi_1$ on the starting board. Let $\Pi$ be the orthogonal projector onto the
row-minus-column mode. At the rook stationary point, set
$G^0=(S_T+\Pi)^{-1}$, $\mathsf G(i,j)=a_i^\top G^0a_j$. The stationary class requires
$\max p_e\le C/k$ and $\lambda_{\min}(S_T+\Pi)\ge\gamma_A>0$. The scaling class above implies
these margins for large $k$. For a pair $E=\{e,f\}$ in distinct rows and columns, define
\[
 \pi^T=w_ew_f e^{L_T(M_{ef})-L_T(M)},\quad
 \pi_1=p_ep_f\{1-\mathsf G(e,f)+\hat c_e+\hat c_f\},\quad
 X=p_ep_f(p_e+p_f)\mathsf G(e,f).
\]
On the binary mask, the ratio $\pi^{\rm ex}/\pi^T$ is exactly
$\exp\{g_T(M_{ef})-g_T(M)\}$. Removing two fibers from each family changes their sums by
$O(k^{-1})$. The dual system in the fixed gauge restores the masses with potentials of order
$O(k^{-1})$. The gap decreases by at most $O(k^{-1})$. By the Schur complement and the
resolvent identity, the retained entries of the Green matrices change by $O(k^{-2})$.
Omitting the indices in the normalized contractions gives
$|\mathfrak F(M_{ef})-\mathfrak F(M)|=O(k^{-1})$, where $\mathfrak F$ is evaluated at the
profile of each mask. Consequently $\Delta_E g_T=g_T(M_{ef})-g_T(M)=O(k^{-2})$ by (T).
Two applications of rook insertion and the stability of $p,\hat c$ give $\pi^T=\pi_1+O(k^{-4})$
(Corollary C for rooks). This gives $\pi^{\rm ex}-\pi_1=O(k^{-4})$, which is statement (R0).

\paragraph{The third square [D].}
When we occupy a third square $c$ in another row and another column, the pair remains available
in the binary minor $M_c$. Applying (R0) to both masks would give only $O(k^{-4})$ for the
difference. Compare $M,M_{ef},M_c,(M_c)_{ef}$ in $\Delta_c\Delta_E(\Phi/k)$ using two mass
parameters. Removing $E$ has parameter $s$, removing $c$ has parameter $t$, and the inverse
size is $k^{-1}/(1-2s-t)$. The kernels and inverses defining $\Phi$ are bounded on fixed
complex disks. The expression divided by $1/k$ has a bounded mixed derivative there.
Integrating over $[0,1/k]^2$ gives $|\Delta_c\Delta_E(\Phi/k)|=O(k^{-3})$. The same bound for
the remainder follows from two applications of $(T^+)$. Then, in
\[
 \pi^{\rm ex}-\pi_1=(\pi^T-\pi_1)+\pi^T(\exp(\Delta_E g_T)-1)
\]
the second summand changes by $O(p_ep_f/k^3)$ by the product rule and
$\Delta_c\pi^T=O(p_ep_f/k)$. For the first we retain $p_ep_f$ in the product of the insertions.
The second differences of $\mathsf G,\hat c$ and the change of order $O(k^{-3})$ in $r_e$
in a distant minor give $|\Delta_c(\pi^T-\pi_1)|\le Cp_ep_f/k^3$. This proves (R1) with weight
$\varpi(E)=k^2p_ep_f$. We now average over the square occupied in a row $r$ distinct from the
rows of the pair. Commutation of three minors and $\sum_{c\in r}\pi^{\rm der}_c=1$ give the
consistency defect of $\pi_1$ under this average,
\[
  D_r=\sum_{\substack{c\in r\\ M_c\supset E}}\pi^{\rm der}_c(M)\pi_1(M_c)-\pi_1(M),
       \qquad |D_r|\le C\varpi(E)/k^5.
\]
To prove this, set $R_\pi=\pi^T-\pi_1$. Commutation gives three summands in $D_r$.
The first is $\pi^T(M)\sum_c\pi_c^{\rm der}(M_{ef})(e^{\Delta_E r_c}-1)$.
Stability (S) of the insertion error gives $\Delta_E r_c=O(k^{-3})$.
The second is $-\sum_c\pi_c^{\rm der}(M)\Delta_cR_\pi$ over the squares that retain the pair.
The preceding increment bound applies here. The third,
$R_\pi(M)\sum_{c\text{ lost}}\pi_c^{\rm der}(M)$, involves at most two squares, each of
mass $O(1/k)$. Combining the three bounds gives (R2).

On binary masks the gap and the factorization of the stationary point give a lower bound
depending only on the class for $kp_e$ on the support, so that absolute bounds for a fixed
number of pairs can be converted into weighted bounds. Extending this conversion to arbitrary
forbidden squares or repeating it for growing groups would require another estimate.
Two collections of blocks with a third vertex have $\hat c_e=\hat c_f=\mathsf G(e,f)=0$ and
$\pi^{\rm ex}=\pi_1$ exactly [O]. In these families the error vanishes for any positive value
of the constant in the bound.

\paragraph{A pair and a square [D].}
On the board $B$ that restores the diagonal restriction of a single line $E$, the logarithm of
the count changes by $\log(1-\pi^{\rm ex})$. To subtract the change in the predictor from this,
let $\lambda=L(B)-L_T(M)$. With $p=p_e,q=p_f$, the first expansion gives
$\lambda=-\pi_1+O(k^{-4})$ (Theorem 1). The second (Theorem 2) gives
\[
       \lambda=\log(1-\pi_1)+X+O(k^{-5}).
\]
At the trial point, a rank-two update of the axial operator incorporates the two squares, and
a Schur complement incorporates the line. The Newton displacement has component $pq+O(k^{-3})$
on the four fibers and $2pq+O(k^{-3})$ on the added line.
The first variation and half the second variation of the action, the local factor, and the
trace of the logarithm leave an analytic function of seven variables of order $O(k^{-1})$.
The degree-four polynomial for $\lambda+\pi_1$ is
$-p^2q^2/2+pq(p+q)\mathsf G(e,f)$ [O, rational coefficients].
The denominators at the origin equal one and the Taylor bound is uniform. Thus the error
$g^\sharp=\log(1-\pi^{\rm ex})-\lambda$ retains the difference between the two logarithms and
$-X$ (Corollary 5).

The average over pairs already includes the effect of each diagonal separately. If all four
squares are available, it remains to measure what imposing both restrictions at once adds. Let
$E'$ be the other pair and $F_\square=E\cup E'$. The update on four squares and the two-dimensional
Schur complement~\cite{hj} give a function of eighteen small arguments.
After subtracting the two single-line expressions taken on the same mask, its degree-four
polynomial is $2\prod_{i\in F_\square}p_i$ [O]. This is the predictor excess (Theorem 3), with
remainder $O(k^{-5})$. The two pairs are incompatible and leave the same minor. If the completion
counts are positive, the exact avoidance excess is $\log[(1-2\pi^{\rm ex})/(1-\pi^{\rm ex})^2]$.
Its first term is $-(\pi^{\rm ex})^2=-\prod_{i\in F_\square}p_i+O(k^{-5})$ by (R0) and by the
factorization of $p_i(1-p_i)$. The joint contribution to the error is therefore
$-3\prod_{i\in F_\square}p_i+O(k^{-5})$ (Corollary 6). Taking the boundary mean counts this
difference once per complete square and gives $-S$ in~\eqref{eq:main}.

